% Modernised reading — fonds Alexandre Grothendieck, shelfmark 162-1, pages 1-9.
%
% This is an INTERPRETATION, not a transcription. In any dispute of substance the
% transcription governs: whoever cites Grothendieck must cite batch-01.fr.tex.
%
% Pass: Opus 5 (claude-opus-5), 2026-08-16 — first pass, unchecked against the
% pages by a human. The folder was read whole; it holds one batch, pages 1-9.
%
% This folder has no mathematics in it, and the reading says so at the outset
% rather than manufacturing some. What it has instead is a list of documents,
% and what a modernised reading can honestly do with a list of documents is
% identify the items — say under what name each is known now — without asserting
% anything about who read what, or when, or to what effect. Every identification
% below is marked as ours. Where the hand does not settle a name, it is left
% unsettled: a lending register is exactly the kind of document where a
% confidently wrong identification would be repeated.
\documentclass[11pt,a4paper]{article}
% Le préambule commun à toutes les transcriptions du fonds Grothendieck.
%
% Il tient en une idée : **l'incertitude se compose, elle ne se résout pas.**
% Une transcription qui lisse un mot douteux a détruit la seule chose pour
% laquelle on la faisait — un lecteur doit pouvoir savoir, à chaque endroit,
% ce qui était sur le feuillet et ce qui a été ajouté. Les six macros qui
% suivent sont donc l'appareil critique, et non de la décoration.
%
% Les mêmes commandes sont reconnues par `scripts/render.mjs`, qui produit la
% vue de lecture du site. Une macro ajoutée ici doit l'être aussi là-bas,
% faute de quoi le rendu échouera bruyamment — ce qui est voulu : un rendu
% silencieusement faux est pire qu'un rendu absent.

\ProvidesPackage{grothendieck}[2026/08/08 Transcriptions du fonds Grothendieck]

\usepackage{amsmath,amssymb,amsthm}
\usepackage{mathrsfs}
\usepackage{stmaryrd}
% Les diagrammes commutatifs sont partout dans ce fonds : c'est la langue
% même de la géométrie algébrique de Grothendieck, et une transcription qui
% les rendrait en prose ne transcrirait rien.
\usepackage{tikz-cd}
% Français par défaut — c'est la langue des feuillets — mais l'anglais est
% chargé aussi : la traduction et le résumé sont des documents anglais, et la
% typographie française (espaces avant les deux-points) y serait une faute.
% Ces documents basculent par \selectlanguage{english} en tête de corps.
\usepackage[english,french]{babel}
\usepackage{fontspec}
\usepackage{geometry}
\geometry{a4paper,margin=25mm,marginparwidth=28mm}
\usepackage{marginnote}
\usepackage{xcolor}
% ulem, not soul. `soul`'s \st and \ul are built on a character-by-character
% scan that XeLaTeX's font machinery breaks: the document fails with an
% undefined control sequence rather than degrading. ulem does the same two jobs
% and is Unicode-engine safe. `normalem` stops it hijacking \emph, which the
% transcriptions use for Grothendieck's own emphasis.
\usepackage[normalem]{ulem}
\usepackage{hyperref}
\hypersetup{colorlinks=true,linkcolor=black,urlcolor=black,pdfborder={0 0 0}}

% --- Métadonnées du lot -----------------------------------------------------
% Reprises telles quelles par le script de rendu, qui les affiche en tête de
% la vue de lecture. Elles fixent ce que le fichier prétend couvrir : sans
% elles, une transcription flotte sans point d'attache dans un fonds de seize
% mille feuillets.

\newcommand{\folder}[1]{\gdef\grothFolder{#1}}
\newcommand{\batch}[1]{\gdef\grothBatch{#1}}
\newcommand{\leaves}[2]{\gdef\grothFirst{#1}\gdef\grothLast{#2}}
\newcommand{\foldertitle}[1]{\gdef\grothFoldertitle{#1}}
\gdef\grothFolder{?}\gdef\grothBatch{?}\gdef\grothFirst{?}\gdef\grothLast{?}\gdef\grothFoldertitle{}

% --- Le résumé --------------------------------------------------------------
%
% Il ouvre la lecture modernisée, sous le titre « Résumé », et il est composé
% en petit. Ce n'est pas un ornement : c'est le seul passage du document qui
% ne soit pas des mathématiques, et un lecteur doit pouvoir le reconnaître
% avant de l'avoir lu — puis le sauter s'il connaît déjà le sujet, ou s'y
% arrêter s'il ne le connaît pas. Le filet qui le suit marque la reprise du
% corps.

\newenvironment{resume}
  {\small\setlength{\parskip}{0.45em}}
  {\par\medskip\hrule\medskip}

% --- Mots-clés ---------------------------------------------------------------
%
% En anglais, en clôture du résumé de la lecture modernisée : le vocabulaire
% moderne sous lequel chercher ce que les feuillets construisent. Le manifeste
% du site (`npm run manifest`) les extrait pour étiqueter la cote entière —
% cette ligne est la source unique des tags, il n'y a pas de fichier à côté.
\newcommand{\keywords}[1]{\par\smallskip\noindent
  {\footnotesize\textsc{Keywords} --- \textit{#1}}\par}

% --- L'appareil critique ----------------------------------------------------

% Le numéro de feuillet, en marge. C'est l'ancre qui permet de régler un
% désaccord sur une formule en regardant *un* feuillet plutôt qu'une chemise
% de six cents. Le site s'en sert aussi pour tourner les pages du fac-similé
% quand on fait défiler la transcription.
\newcommand{\leaf}[1]{%
  \marginnote{\footnotesize\color{gray!70}\textsf{#1}}%
  \label{leaf:#1}\ignorespaces}

% Illisible : on ne devine pas. Un mot inventé qui se lit comme les autres est
% le pire résultat possible.
\newcommand{\ill}{\textcolor{red!60!black}{[\,\ldots\,]}}

% Lecture douteuse : le mot est donné, et signalé comme incertain.
\newcommand{\uncertain}[1]{\uline{#1}}

% Ajout de l'éditeur — une lettre manquante, un mot sous-entendu.
\newcommand{\add}[1]{\textcolor{blue!50!black}{[#1]}}

% Biffé par Grothendieck lui-même. Ce qu'il a rayé fait partie du document :
% une rature dit souvent où la pensée a changé de direction.
\newcommand{\struck}[1]{\sout{#1}}

% Note du transcripteur, sur l'état matériel du feuillet.
\newcommand{\note}[1]{\par\smallskip{\small\color{gray!60!black}\textit{[#1]}}\par\smallskip}

% Note marginale de Grothendieck, distincte du corps du texte.
\newcommand{\marginal}[1]{\par\smallskip\noindent\hspace{1em}%
  {\small\color{gray!60!black}#1}\par\smallskip}

% --- Titre ------------------------------------------------------------------

\AtBeginDocument{%
  \begin{center}
    {\large\bfseries Fonds Alexandre Grothendieck --- cote n\textsuperscript{o}~\grothFolder}\\[2pt]
    {\itshape\grothFoldertitle}\\[4pt]
    {\small feuillets \grothFirst--\grothLast\ (lot \grothBatch)}\\[2pt]
    {\footnotesize\color{gray!60!black}Transcription établie d'après le fac-similé numérisé
     par l'Université de Montpellier.\\
     Les lectures incertaines sont soulignées, les illisibles notées [\,\ldots\,],
     les ajouts entre crochets.}
  \end{center}
  \vspace{1em}}

\endinput


\folder{162-1}
\pages{1}{9}
\dating{[à partir de 1967] — le groupe « Dossiers rassemblés par Grothendieck sur différentes thématiques » (161-1 à 162-6) est daté [après 1961-vers 1977]}
\watermark{Édition de démonstration\\interprétation personnelle de l'œuvre}
\foldertitle{Papiers prêtés [liste de noms et de références documentaires] : notes manuscrites (s.d.) — lecture modernisée du dossier entier}

\begin{document}

\section*{Résumé}

\begin{resume}
Ce dossier ne contient aucune mathématique. Il contient un registre de prêts :
neuf pages où, dans une colonne de gauche, figure un nom, et dans une colonne de
droite, ce que la personne a emporté. La plupart des lignes sont barrées.

Il faut dire tout de suite pourquoi un tel objet mérite d'être lu. Vers la fin
des années soixante, une grande partie de ce qui s'écrivait en géométrie
algébrique n'était pas imprimée. Les séminaires circulaient en polycopiés, les
rédactions passaient de main en main sous forme de notes manuscrites, les
articles s'échangeaient en tirés à part, et les lettres --- de Serre, de
Deligne, d'Atiyah --- avaient valeur de publications. Il n'y avait ni
photocopieuse à chaque étage ni serveur de prépublications. Un exemplaire
existait, et il était quelque part.

Un registre de prêts est donc, à cette date, un instrument de travail
nécessaire : sans lui on ne sait plus qui a le seul exemplaire d'un chapitre. Et
c'est aussi, quarante ans plus tard, une photographie. Il montre, sans intention
de le montrer, ce qui circulait dans un bureau donné à un moment donné : quels
textes, entre quelles mains, et dans quelles combinaisons.

Ce que la photographie donne à voir tient en trois traits. Le premier est
l'omniprésence des séminaires non publiés : SGA 4, SGA 5, SGA 6, SGA 7, EGA IV et
EGA V circulent en fascicules et en exposés isolés, souvent désignés par un
numéro d'exposé sans autre précision. Le deuxième est la place des lettres :
une lettre à Serre, une lettre à Dieudonné, deux lettres de Deligne sur les
connexions intégrables sont prêtées comme le serait un article. Le troisième est
la coexistence, sur la même page, de sujets qu'on ne penserait pas voisins : la
cohomologie cristalline et les espaces de Banach, les catégories tannakiennes et
la topologie différentielle.

Les strates du registre méritent enfin d'être signalées. Les entrées barrées
sont vraisemblablement des retours, et l'écrasante majorité l'est ; mais le
registre ne le dit nulle part, et cette lecture ne l'affirme pas. Deux mentions
marginales --- \emph{redemander} --- disent en revanche clairement ce qu'elles
disent.

\keywords{mathematical correspondence, preprint circulation, Séminaire de Géométrie Algébrique, IHES, provenance}
\end{resume}

\pagerange{1}{2}
\section*{Comment le registre est fait}

La chemise porte deux mots à l'encre --- \emph{Papiers prêtés} --- et, au
crayon, un chiffre romain qui se lit III : ce dossier serait donc le troisième
d'une série dont le reste n'est pas ici.\footnote{Le chiffre est incertain et la
transcription le marque comme tel ; une indication entre parenthèses le suit,
restée illisible. On ne conclut pas.}

Chaque page se lit en deux colonnes. À gauche l'emprunteur, souvent souligné,
et parfois suivi d'un chiffre --- 1 ou 2 --- qui revient trop régulièrement pour
être un nombre de pièces.\footnote{Il compte vraisemblablement les rappels, mais
le registre ne l'explique jamais et l'on ne le suppose pas. Un nom peut
d'ailleurs reparaître à plusieurs pages avec un chiffre différent.} À droite,
ce qui a été emporté, en une ou plusieurs lignes.

Une convention gouverne toute la lecture, et elle est de nous. Environ les deux
tiers des entrées sont barrés. Sur un registre de prêts, la lecture naturelle
est qu'une entrée barrée est une entrée soldée --- le document est revenu. Rien
sur la page ne l'affirme, et il reste concevable qu'une rature signale au
contraire un prêt annulé ou une inscription erronée. La transcription conserve
donc les ratures comme ratures, et cette lecture ne les traduit pas.\footnote{La
proportion de passages illisibles est ici bien plus élevée que dans un dossier
mathématique : les ratures sont lourdes, plusieurs se recouvrent, et un nom rayé
deux fois a généralement disparu. La transcription compte une soixantaine de
lacunes signalées. Aucune n'a été comblée, ici pas plus que là.}

\pagerange{2}{9}
\section*{Ce que le registre permet d'identifier}

Ce qui suit relève un petit nombre d'entrées dont la désignation, si abrégée
soit-elle, suffit à identifier l'objet sans risque. C'est le seul travail
proprement « modernisant » qu'un tel document autorise.\footnote{Toutes les
identifications de cette section sont les nôtres. Aucune ne figure sur la page
autrement que sous la forme abrégée reproduite ; aucune n'est certaine au sens
où le serait une référence complète ; et aucune ne dit quoi que ce soit de
l'usage qui a été fait du document prêté.}

\subsection*{Cohomologie $p$-adique et cristalline}

Le registre porte, à deux endroits différents et pour deux emprunteurs
différents, une entrée \emph{Monsky--Washnitzer I, II}. Le nom désigne la
cohomologie construite par ces deux auteurs pour les variétés affines lisses en
caractéristique $p$, publiée en deux parties.\footnote{Le nom est écrit très
vite et la seconde moitié en est peu formée ; la lecture retenue est la nôtre.
Ce qui la soutient est le contexte immédiat --- l'entrée voisine, chez le même
emprunteur, est « Torsion en cohomologie cristalline » --- et non la seule
graphie.}

Dans le même voisinage on relève : \emph{Torsion en cohomologie cristalline},
sous le nom de Grothendieck lui-même ; \emph{notes sur groupes de Barsotti--Tate
et cristaux} ; \emph{notes Cartier sur vecteurs de Witt} ; et, chez Berthelot,
plusieurs papiers de Katz --- le théorème de Dwork, la théorie des points
singuliers réguliers, et les connexions nilpotentes avec Gauss--Manin, en deux
versions.

\subsection*{Complexe cotangent et déformations}

Chez Illusie, une ligne : \emph{Complexe cotangent relatif de Quillen}. Elle
voisine avec des notes sur les extensions infinitésimales, une lettre de
Quillen, et une entrée sur l'extension d'Atiyah pour $\Omega^{1}$. Le groupe
est cohérent : c'est la théorie des déformations dans l'état où elle se trouvait
avant sa rédaction définitive.

Sur une autre page, chez un emprunteur dont le nom est rayé,
\emph{Schlessinger, Functors of Artin rings} --- l'article qui donne les
critères de pro-représentabilité d'un foncteur de déformations. Il reparaît plus
loin sous la forme \emph{Schlessinger, thèse}.

\subsection*{Catégories tannakiennes}

Le nom de Saavedra revient quatre fois, et à chaque fois avec des entrées de
même famille : \emph{notes sur polarisations dans catégories tannakiennes},
\emph{notes sur structure de Hodge}, une lettre de Mumford sur la description
transcendante des motifs en caractéristique 0, \emph{note filtrations et groupes
paraboliques}, et à deux reprises des notes sur les catégories de
Picard.\footnote{Le rapprochement avec le sujet 22 du dossier 66 --- « structure
des catégories abéliennes avec produit tensoriel rigide \ldots\ préliminaire
algébrique à la théorie des motifs », annoté \emph{travail Saavedra en train} ---
s'impose et vaut d'être noté, les deux dossiers étant du même fonds. Il ne
prouve rien : deux documents peuvent nommer la même personne sans se rapporter
au même moment.}

\subsection*{Conjectures standard et cycles}

Chez B.~Mazur : \emph{notes Kleiman sur standard}. L'abréviation est celle des
conjectures standard sur les cycles algébriques.\footnote{Le mot est incertain
et la transcription le marque. C'est l'une des identifications les plus
vraisemblables du registre --- Kleiman est l'auteur du texte de référence sur
ces conjectures --- et l'une de celles qu'il faut le plus se garder d'affirmer,
puisque le seul appui est un mot mal formé.} Sur la même page,
\emph{correspondance avec Griffiths}, barrée.

Ailleurs : \emph{Matsusaka}, avec trois titres dont un critère d'équivalence
algébrique ; \emph{Manin, variétés rationnelles} ; \emph{article Mumford,
cycles}.

\subsection*{Séminaires et fascicules}

Le registre est enfin, et surtout, un registre de séminaires. On y relève SGA 4
fascicule 1 et les exposés IV et IX ; SGA 5 exposés IV, XII (§§ 9 et 10) et
XVIII ; SGA 6 exposés I, III et XIV ; SGA 7 exposés I à V, VI, VII, XX et XXI ;
EGA IV chapitres 16 à 19 et 20 à 21 ; EGA V.

La régularité de ces désignations mérite d'être relevée. On ne prête pas « le
séminaire » : on prête un exposé, désigné par son numéro romain, comme on
prêterait un chapitre. C'est le signe que ces textes circulaient à l'unité, et
qu'un numéro d'exposé suffisait entre gens du même bureau à désigner un objet
sans ambiguïté.

\subsection*{Ce qui ne relève pas de la géométrie algébrique}

Deux entrées font exception et méritent d'être signalées pour cette raison.
Chez Thomas : \emph{papiers Lindenstrauss, Pełczyński}, et un \emph{survey} de
Lindenstrauss --- c'est-à-dire de la théorie des espaces de Banach. Chez
Poénaru : \emph{Atiyah--Hirzebruch, über Cohomologie-Operationen}, et une revue
du livre de Bishop par Stolzenberg. Et chez un emprunteur dont le nom n'est pas
assuré : \emph{produits tensoriels topologiques (Grothendieck)}, qui renvoie à
son propre travail d'analyse fonctionnelle, antérieur de quinze ans à tout le
reste du registre.\footnote{C'est la seule entrée du dossier qui renvoie à la
première période de son \oe uvre. Elle est notée sans commentaire, comme
n'importe quelle autre.}

\pagerange{1}{9}
\section*{Ce que ce dossier peut et ne peut pas servir à établir}

Il vaut la peine de dire ce qu'un tel document autorise, parce que la tentation
d'en tirer davantage est réelle.

Il établit qu'un texte existait sous une forme prêtable à une date antérieure à
celle du registre, et qu'il portait, dans l'usage courant de ce bureau, la
désignation abrégée qui figure sur la page. C'est déjà beaucoup pour un corpus
dont une grande part n'a jamais été imprimée.

Il n'établit rien d'autre. Il ne dit pas que l'emprunteur a lu ce qu'il a
emporté, ni ce qu'il en a fait, ni que le prêt a précédé un travail plutôt que
l'inverse. Il ne date rien avec précision : la seule date portée sur les neuf
pages est un 10.I.66 accolé à une entrée, et une lettre du 21.9.67 mentionnée
sur une autre.\footnote{La première est antérieure au « [à partir de 1967] »
retenu par l'inventaire, la seconde le confirme. Rien dans le dossier n'explique
comment l'inventaire est arrivé à sa borne, et l'écart n'est pas résolu ici.}

Et il ne dit rien sur la paternité de quoi que ce soit. Un registre de prêts
ressemble par sa forme à une bibliographie, mais il n'en est pas une : il
enregistre un déplacement d'objets, non une dette intellectuelle. Lire les deux
comme équivalents serait la seule erreur grave qu'on puisse commettre sur ce
dossier.

\end{document}
